% Prüfungsprogramm Aikido Shinki Rengo · Daitoryu · Itto-den Shinki Toho
%
% Bauen: ./build.sh (Heft) bzw. ./build.sh entwurf (mit Fragen-Markern).
% Alle Inhalte kommen aus mitori-daten/daten/*.yaml (über tools/daten.py); gesetzt von lua/mitori-heft.lua.
% build.sh setzt \mitoridaten (Submodul, sonst MITORI_DATEN_DIR).
\providecommand{\mitoridaten}{../mitori-daten}
\documentclass[10pt]{article}
\usepackage{stil/shinki}

\begin{document}

% Titel ------------------------------------------------------------------------
\setzerubrik{}{}
\vspace*{18mm}
\begin{center}
  \includegraphics[height=46mm]{grafik/shinki-logo.pdf}\par\vspace{16mm}
  {\sffamily\bfseries\fontsize{34}{38}\selectfont Prüfungsprogramm}\par\vspace{4mm}
  {\color{shu}\rule{40mm}{1.4pt}}\par\vspace{8mm}
  {\sffamily\fontsize{14}{20}\selectfont Aikido Shinki Rengo\\Daitoryu Aiki Jujutsu Bokuyokan Europa\\Itto-den Shinki Toho}\par\vspace{10mm}
  {\jpschrift\fontsize{16}{20}\selectfont\color{usuzumi}合気道　大東流合気柔術　一刀伝神氣刀法}\par
  \vfill
  {\sffamily\fontsize{9}{12}\selectfont\color{usuzumi}Shinkiryu · Stand 2026 · Daten \directlua{mitori.stand()}\ifentwurf\\\textcolor{shu}{Entwurf mit offenen Fragen}\fi}
\end{center}
\clearpage

% Inhalt -----------------------------------------------------------------------
\hochformat\setzerubrik{Prüfungsprogramm}{Inhalt}
{\sffamily\bfseries\fontsize{20}{24}\selectfont Inhalt}\par\vspace{2mm}{\color{shu}\rule{30mm}{1.2pt}}\par\vspace{6mm}
{\makeatletter\renewcommand\contentsname{}\let\section\@gobbletwo
 \renewcommand*\l@section[2]{\addvspace{3mm}{\sffamily\bfseries\fontsize{11}{14}\selectfont\color{shu}#1\hfill\color{sumi}#2}\par}
 \renewcommand*\l@subsection[2]{{\fontsize{10}{13}\selectfont\hspace*{5mm}#1\dotfill#2}\par}
 \@starttoc{toc}\makeatother}

\textseite{Prüfungsprogramm}{Hinweise}{Zum Gebrauch}
% Hinweise zum Gebrauch des Hefts (vorne im Heft).
{\fontsize{10.5}{14.5}\selectfont\setlength\parskip{2.5mm}
Das Heft zeigt jedes Programm in zwei Ansichten, beide aus denselben Daten gesetzt:

Vorne in jedem Teil steht eine Übersicht der Grade mit ihren Gruppen und Wartezeiten.

\textbf{Ein Grad je Seite.} Alles, was zu einem Grad gehört, nach Gruppen und Angriffen
geordnet. Gedacht zum Lernen und Mitnehmen.

\textbf{Matrix.} Angriff mal Grad im Querformat, wie die bisherigen Prüfungsordnungen.
Gedacht für Prüferinnen und Prüfer und für Stichproben über mehrere Grade.

\textbf{Schreibweise.} Namen in Umschrift; \textit{kursiv} steht die Ausführung oder Variante.
Längenzeichen (ô, û) wie in der jeweiligen Vorlage: Aikido ohne, Daitoryu und Shinki Toho mit.

\textbf{Japanisch.} Aikido: aus den Technik-Listen unseres Sensei. Shinki Toho: Grade und
Kamae von shinkiryu.de, Formen 1–21 von den Titelkarten der Prüfungsvideos.
Daitoryu: Die Vorlagen enthalten keine Schriftzeichen; die Angaben in \textcolor{usuzumiv}{hellem Grau}
sind ein Vorschlag und müssen noch von einem Lehrer bestätigt werden.

\textbf{Vorlagen.} Aikido Shinki Rengo, Prüfungsordnung, Stand 01.01.2022.
Daitoryu Aiki Jujutsu Bokuyokan Europa, Prüfungsordnung Kyu- und Dan-Grade, 01.04.2008.
Itto-den Shinki Toho: Grade und Prüfungsberechtigung, Kamae-Tabelle, Suburi,
Ô-dachi-Tabelle (Stand 25.08.2010) und Ô-dachi-Kommentare.
\ifentwurf

\textbf{Entwurf.} Rote Marken wie \frage{A-K-02} verweisen auf eine offene Frage in
\texttt{KORREKTUREN.md}. Im fertigen Heft fehlen sie.\fi
\par}


% Aikido -----------------------------------------------------------------------
\teil{Aikido Shinki Rengo}{合気道}{Prüfungsprogramm Kyu- und Dan-Grade\\Stand 01.01.2022 · überarbeitet 2026}
\textseite{Aikido Shinki Rengo · Prüfungsprogramm}{Vorbemerkung}{Vorbemerkung}
{\fontsize{10.5}{14.5}\selectfont\setlength\parskip{2.5mm}\input{\mitoridaten/daten/aikido-vorbemerkung.tex}\par}
\inhalt{Grade \& Programm}\ansicht{grade-aikido}
\inhalt{Kyu-Grade: ein Grad je Seite}\ansicht{tusche-aikido-1}
\inhalt{Dan-Grade: ein Grad je Seite}\ansicht{tusche-aikido-2}
\inhalt{Kyu-Grade: Matrix}\ansicht{matrix-aikido-1}
\inhalt{Dan-Grade: Matrix}\ansicht{matrix-aikido-2}

% Daitoryu ---------------------------------------------------------------------
\teil{Daitoryu Aiki Jujutsu}{大東流合気柔術}{Bokuyokan Europa\\Prüfungsprogramm Kyu- und Dan-Grade\\Stand 01.04.2008 · überarbeitet 2026}
\inhalt{Grade \& Programm}\ansicht{grade-daitoryu}
\inhalt{Kyu-Grade: ein Grad je Seite}\ansicht{tusche-daitoryu-1}
\inhalt{Dan-Grade: ein Grad je Seite}\ansicht{tusche-daitoryu-2}
\inhalt{Kyu-Grade: Matrix}\ansicht{matrix-daitoryu-1}
\inhalt{Dan-Grade: Matrix}\ansicht{matrix-daitoryu-2}

% Shinki Toho ------------------------------------------------------------------
\teil{Itto-den Shinki Toho}{一刀伝 神氣刀法}{Grade, Kamae, Suburi und Ô-dachi-Formen\\Stand 2026}
\inhalt{Grade, Zulassung, Kamae und Suburi}\ansicht{shinkitoho-grade}
\inhalt{Ô-dachi je Stufe}\ansicht{shinkitoho-odachi}
\inhalt{Ô-dachi: alle Formen}\ansicht{shinkitoho-matrix}

\end{document}
